% problem-000217 10 乙酸正丁酯实验合成
\section*{10. 乙酸正丁酯实验合成}
⼄酸正丁酯是优良的有机溶剂，在⽯油加⼯、制药⼯程等领域均有⼴泛应⽤。在实验室中，⼄酸正丁酯可

通过正丁醇和冰醋酸在强酸催化及加热条件下合成，实验装置如右图所⽰。反应中⽣成的⽔可以与正丁醇或⼄酸正丁酯形成共沸混合物被蒸出。使⽤分⽔器可在反应进⾏的同时分离出⽣成的⽔，并使有机相流回反应瓶中。实验步骤如下：（1）如右图所⽰装配仪器。

\subsection*{(2)}
关闭分⽔器旋塞并在分⽔器带旋塞⼀侧预先加⼊适量的 ① ，再向反应瓶中加⼊5.0 mL (55 mmol)正丁醇、 ② 和 ③ ，加⼊1 粒沸⽯。

\subsection*{(3)}
将反应液温度控制在 ④ 并保持15分钟，然后升温⾄回流状态进⾏分⽔。当反应不再有⽔⽣成时，停⽌加热。合并烧瓶和分⽔器中的液体，⽤分液漏⽃将⽔层分出。依次⽤5 mL 10%碳酸钠⽔溶液及5 mL⽔洗涤有机层。

（图：）

\subsection*{(4)}
有机层⽤⽆⽔硫酸镁⼲燥。常压蒸馏，收集 $1 2 4 { \sim } 1 2 6 ^ { \circ } \mathrm { C }$ 的馏分。

各化合物及相关共沸物常压下的沸点（单位：°C）

\textit{（原卷此处含表格/仪器试剂表，详见 PDF 或 MD 源文件。）}

\subsection*{10-1}
根据上述实验流程描述，写出图中所⽰实验装置中的三处错误。

\subsection*{10-2}
在上述编号处分别填⼊最合适的选项：①处：A. 正丁醇 B. 冰醋酸 C. ⽔

②处：A. 25 mmol 冰醋酸 B. 50 mmol 冰醋酸 C. 60 mmol 冰醋酸③处：A. 0.05 mL 浓硫酸 B. 1 mL 浓硫酸 C. 0.05 mL 稀硫酸 D. 1 mL 稀硫酸④处：A. $0 ~ ^ { \circ } \mathrm { C }$ B. $2 5 ^ { \circ } \mathrm { C }$ C. $8 0 ~ ^ { \circ } \mathrm { C }$

\subsection*{10-3}
解释步骤（3）中为何要⽤分⽔器将⽔分出。
